\documentclass[12pt]{article}
\usepackage[utf8]{inputenc}
\usepackage[T1]{fontenc}
\usepackage{hyperref}
\title{Tombelala Language Resources for K-12 Curriculum Development}
\author{Compiled by Claude Sonnet 5 (claude-sonnet-5) for Andr\'as Kornai}
\date{}
\begin{document}
\maketitle

\section{Introduction}
Tombelala (ISO 639-3: ttp; Glottocode tomb1245) is an Austronesian language of Central Sulawesi, Indonesia. It is also called Mbelala, Belala or Baria. Wikipedia places it in the Celebic branch, related to Pamona. Its speakers live mainly in Lamberea village, Morowali Regency. Wikipedia cites about 1,100 speakers (1991). Joshua Project gives much higher figures for the wider Tombelala people. Joshua Project also reports that children do not use the language and that it is not taught in schools.

\section{Primary Educational Resources}
The searches found no primers, textbooks or school programmes for Tombelala.

\section{Linguistic and Reference Materials}
\begin{itemize}
\item Wikipedia article on Mbelala: \url{https://en.wikipedia.org/wiki/Mbelala_language}
\item Endangered Languages Project entry, based on Ethnologue 2009 and noting 66--76\% lexical similarity with Pamona dialects: \url{https://endangeredlanguages.com/lang/2797}
\item Glottolog reference: Mead, David and Mead, Melanie 1991. Only the reference title was seen in the results, so its content is unverified: \url{https://glottolog.org/resource/reference/id/96623}
\item Verbix endangered-languages page: \url{https://docs.verbix.com/EndangeredLanguages/Tombelala}
\end{itemize}

\section{Religious and Community Resources}
\begin{itemize}
\item Joshua Project profile of the Tombelala people: \url{https://joshuaproject.net/people_groups/15548}
\item Asia Harvest people-group profile (PDF): \url{https://people-groups.asiaharvest.org/Indonesia/Tombelala.pdf}
\item Global Recordings Network page for ttp: \url{https://globalrecordings.net/en/language/ttp}
\end{itemize}
The result snippets did not show whether any scripture or audio exists.

\section{Digital and Audio Resources}
No audio, video or corpus resources were confirmed.

\section{Curriculum Development Context}
Tombelala is threatened, and young people are not acquiring it fully. No teaching materials were found. Any curriculum would need to begin with documentation and with teaching it to children, working with the community.
\end{document}